\chapter{The Algebraic Trap and Topological Excitations}

Chapter 3 identified the thermodynamic endpoint of holographic scaling as a nilpotent zero-mode sector.  We now describe the topology of the trap that catches these modes.  The central observation is that the boundary possesses fixed loci: spatial fixed lines produced by nonsymmorphic glide/Klein geometry and spectral fixed lines produced by scale reflection.  These two fixed sets are formally parallel:
\begin{equation}
    k_2=0
    \qquad\longleftrightarrow\qquad
    \operatorname{Re}(s)=\frac12.
\end{equation}
The first is the reciprocal-space fixed line of a Brillouin Klein bottle; the second is the critical line of the Riemann scale duality.  The chapter develops this spatial--spectral dictionary and connects it to Dirac resolvents, Fourier--Mellin transforms, Möbius parity, and Witten index stability.

\section{The Dirac--Fourier--Mellin Factorization}

The algebraic bridge from bulk differential geometry to boundary spectral theory is the Fourier--Mellin transform of the Dirac operator.  In the Pauli-biquaternion representation, the relevant resolvent takes the schematic form
\begin{equation}
    D_{FM}(s,X)=-\sigma_3(sI_2-X),
\end{equation}
where $X$ is a biquaternionic matrix coordinate and $s$ is the Mellin spectral parameter.

The determinant of the resolvent detects poles and zero modes.  When the determinant vanishes, a boundary state is trapped:
\begin{equation}
    \det(sI_2-X)=0.
\end{equation}
In the nilpotent boundary limit, this condition collapses to the zero-momentum pole
\begin{equation}
    s=0,
\end{equation}
matching the square-zero sector discussed in Chapter 3.

The corresponding formal components are verified in
\begin{center}
\texttt{DiracResolventZeroModeTripotent.lean},\quad
\texttt{KANFourierMellinDirac.lean},\quad
\texttt{DiracFourierMellin.lean}.
\end{center}

\section{The Brillouin Klein Bottle}

A nonsymmorphic glide identifies momentum points by combining reflection with half-translation.  In reciprocal coordinates, the simplest fixed-line model is
\begin{equation}
    (k_1,k_2)\mapsto(k_1,-k_2).
\end{equation}
The fixed locus is therefore
\begin{equation}
    k_2=0.
\end{equation}
When this reflection is combined with reciprocal translations and orientation reversal, the Brillouin zone acquires Klein-bottle topology.  The boundary is nonorientable; excitations transported around the glide cycle return with twisted orientation.

This formalizes the geometric side of the trap:
\begin{equation}
    \text{glide reflection}
    \quad\Longrightarrow\quad
    \text{Klein fixed line}
    \quad\Longrightarrow\quad
    \text{protected boundary modes}.
\end{equation}

Relevant Lean modules include
\begin{center}
\texttt{BrillouinKleinNilpotentAttractor.lean},\quad
\texttt{FixedLineRiemannKlein.lean},\quad
\texttt{GlideDiracSelectionRule.lean}.
\end{center}

\section{The Riemann Scale Reflection}

The spectral analogue of the Klein reflection is the Riemann scale duality.  The completed zeta functional equation relates $s$ and $1-s$.  On the real coordinate, the reflection is
\begin{equation}
    x\mapsto 1-x.
\end{equation}
Its fixed point is
\begin{equation}
    x=\frac12.
\end{equation}
Thus the spectral fixed line is
\begin{equation}
    \operatorname{Re}(s)=\frac12.
\end{equation}

The Lean master certificate records this algebraic kernel as
\begin{equation}
    \texttt{riemannReflectReal}(1/2)=1/2.
\end{equation}
The Klein fixed coordinate is recorded as
\begin{equation}
    \texttt{kleinReflect}(0)=0.
\end{equation}
The affine recentering
\begin{equation}
    x\mapsto x-\frac12
\end{equation}
then sends the Riemann fixed coordinate to the Klein fixed coordinate:
\begin{equation}
    \frac12\mapsto0.
\end{equation}
This is formalized in \texttt{ThesisMaster.lean} as \texttt{riemann\_to\_klein\_fixed\_coordinate}.

\section{Spatial--Spectral Duality}

The dictionary is therefore:
\begin{center}
\begin{tabular}{c|c}
\textbf{Spatial/Klein side} & \textbf{Spectral/Riemann side} \\
\hline
reciprocal coordinate $k_2$ & real spectral coordinate $\operatorname{Re}(s)-1/2$ \\
glide reflection $k_2\mapsto-k_2$ & scale reflection $s\mapsto1-s$ \\
fixed line $k_2=0$ & critical line $\operatorname{Re}(s)=1/2$ \\
Klein nonorientability & functional-equation duality \\
protected nodal modes & spectral zero-mode constraint
\end{tabular}
\end{center}

This is not claimed as a proof of the Riemann Hypothesis.  Rather, it is a theorem-honest fixed-line equivalence: both structures are governed by the same $\mathbb{Z}_2$ reflection geometry after recentering.

\section{Möbius Parity and Witten Index}

The arithmetic parity layer enters through the Möbius function $\mu(n)$.  It behaves as a number-theoretic analogue of fermion parity:
\begin{equation}
    (-1)^F
    \qquad\leftrightarrow\qquad
    \mu(n).
\end{equation}
Square-free integers carry parity according to the number of prime factors; nonsquare-free integers vanish.  Thus the Möbius function naturally combines parity and nilpotent annihilation:
\begin{equation}
    \mu(n)=0
    \quad\text{when}\quad
    n\text{ contains a repeated prime factor}.
\end{equation}

The twisted Witten index is then the protected thermodynamic quantity:
\begin{equation}
    \mathcal{I}(\beta)=\operatorname{Tr}\left((-1)^F e^{-\beta H}\right).
\end{equation}
When supersymmetry is unbroken, this index is independent of $\beta$.  In the boundary architecture, this means thermal deformation does not destroy the protected topological count of zero modes.

The corresponding formal modules include
\begin{center}
\texttt{MobiusWittenKleinIndex.lean},\quad
\texttt{MobiusWittenIndex.lean},\quad
\texttt{SuperchargeSquare.lean}.
\end{center}

\section{Glide Selection Rules and Nodal Protection}

The nonsymmorphic glide imposes selection rules on boundary Dirac modes.  On glide-invariant lines, odd parity sectors can be extinguished while even sectors survive.  This mechanism explains why certain superconducting wallpaper-fermion pairings exhibit protected point or line nodes.

The same algebraic pattern recurs:
\begin{equation}
    G^2=T,
    \qquad
    Q^2=H,
    \qquad
    q^2=0.
\end{equation}
A bulk glide or supercharge squares to a nontrivial even translation/Hamiltonian; at the boundary it degenerates into a nilpotent annihilator.  The trap is therefore simultaneously geometric, spectral, and supersymmetric.

\section{The Non-Hermitian Nielsen--Ninomiya Loophole}

The existence of unpaired chiral zero modes is traditionally obstructed by the Nielsen--Ninomiya theorem, which forces lattice fermions on an orientable Brillouin torus to appear in compensating pairs.  In the present architecture this obstruction is bypassed in two simultaneous ways: the boundary excitations are non-Hermitian exceptional points, whose invariants are non-Abelian braids, and the Brillouin parameter space is not an ordinary torus but a glide-identified Klein bottle.

On the torus, the total exceptional charge is constrained by the braid commutator
\begin{equation}
    B_{\mathrm{total}}=B_xB_yB_x^{-1}B_y^{-1}.
\end{equation}
If the charges commute, this expression is trivial and the ordinary doubling intuition is recovered.  But braid charges need not commute.  This is the formal content of \texttt{ExceptionalBraidTopology.lean}:
\begin{equation}
    \operatorname{Commute}(B_x,B_y)
    \Rightarrow
    B_{\mathrm{total}}=1,
\end{equation}
whereas
\begin{equation}
    \neg\operatorname{Commute}(B_x,B_y)
    \Rightarrow
    B_{\mathrm{total}}\neq1.
\end{equation}

On the Klein bottle the situation is sharper.  The fundamental glide relation is
\begin{equation}
    g a g^{-1}=a^{-1},
\end{equation}
or equivalently
\begin{equation}
    g a g^{-1}a=1.
\end{equation}
Thus the topological partner required for cancellation is not an independent doubler.  It is the glide-reflected inverse of the same exceptional point.  This is formalized in \texttt{ExceptionalKleinGlideEP.lean} by the theorem
\begin{equation}
    \texttt{klein\_word\_trivial\_of\_glide\_inverse}.
\end{equation}
The boundary fermion therefore does not double; it glides.

\section{Möbius--Majorana Self-Pairing}

This glide identification is the non-Hermitian analogue of the Majorana condition.  In the Hermitian setting, a Majorana mode is its own charge conjugate.  On the Brillouin Klein bottle, an exceptional point is its own orientation-reversed glide partner:
\begin{equation}
    a\xrightarrow{\;g\;}gag^{-1}=a^{-1}.
\end{equation}
The anomaly-canceling partner exists, but it is not an additional bulk degree of freedom.  It is the same state after transport around the Möbius cycle.

At the glide fixed line, an even stronger collapse occurs.  If the charge is fixed by the glide,
\begin{equation}
    gag^{-1}=a,
\end{equation}
and the Klein relation also holds,
\begin{equation}
    gag^{-1}=a^{-1},
\end{equation}
then
\begin{equation}
    a^2=1.
\end{equation}
This is proved in Lean as
\begin{equation}
    \texttt{square\_one\_of\_glide\_fixed\_and\_inverse}.
\end{equation}
Physically, the non-Abelian braid charge collapses on the fixed line to a $\mathbb{Z}_2$ involution.  This is the same square-root logic that appears in the projective glide and supercharge laws:
\begin{equation}
    G^2=T,
    \qquad
    Q^2=H,
    \qquad
    q^2=0.
\end{equation}
The bulk square-root symmetry survives as a glide involution, while the boundary degeneration becomes nilpotent.

The resulting slogan is
\begin{equation}
    \boxed{
    \text{EPs do not double on the Klein bottle; they glide.}
    }
\end{equation}

\section{Chapter Summary}

Chapter 4 identifies the topological trap governing boundary excitations.  The Brillouin Klein bottle supplies a spatial fixed line $k_2=0$; the Riemann scale reflection supplies a spectral fixed line $\operatorname{Re}(s)=1/2$; and an affine recentering identifies their algebraic fixed-coordinate structures.  Möbius parity and the Witten index then stabilize the trapped zero modes against thermal deformation.

The chapter's central principle is
\begin{equation}
    \boxed{
    k_2=0
    \quad\longleftrightarrow\quad
    \operatorname{Re}(s)=\frac12
    }
\end{equation}

Chapter 5 will absorb all previous structures into the universal polynomial-operator framework, where rotations, glides, supercharges, tripotents, nilpotents, and Clifford anomaly cancellation become instances of one algebraic language.
